\chapter{Barycentre et lignes de niveau}

\noindent\textit{Le barycentre est le « point d'équilibre » d'un système de points affectés
de coefficients : il généralise le milieu et le centre de gravité. Outil central de la série C,
il permet de réduire des sommes vectorielles, de prouver des alignements ou des concours de
droites, et — couplé à la formule de Leibniz — de décrire de belles \textbf{lignes de niveau}
(cercles, médiatrices, droites).}
\vspace{8pt}

% ═══════════════════════════════════════════════════════════════
\section{Barycentre de points pondérés}

\begin{definition}[Point pondéré et barycentre de deux points]
Un \textbf{point pondéré} est un couple $(A,\alpha)$ où $A$ est un point et $\alpha\in\R$ son \textbf{coefficient} (ou poids). Soient $(A,\alpha)$ et $(B,\beta)$ deux points pondérés tels que $\alpha+\beta\neq0$. Il existe un \textbf{unique} point $G$, appelé \textbf{barycentre} de $\{(A,\alpha),(B,\beta)\}$, vérifiant
\[
\alpha\,\overrightarrow{GA}+\beta\,\overrightarrow{GB}=\vec 0.
\]
\end{definition}

\begin{theoreme}[Caractérisation vectorielle (point de réduction)]
Si $\alpha+\beta\neq0$, le barycentre $G$ de $\{(A,\alpha),(B,\beta)\}$ est l'unique point tel que, pour \textbf{tout} point $M$ du plan,
\[
\alpha\,\overrightarrow{MA}+\beta\,\overrightarrow{MB}=(\alpha+\beta)\,\overrightarrow{MG}.
\]
En particulier (avec $M=A$) : $\ \overrightarrow{AG}=\dfrac{\beta}{\alpha+\beta}\,\overrightarrow{AB}$.
\end{theoreme}

\begin{definition}[Barycentre de $n$ points — isobarycentre]
Soit $\{(A_1,\alpha_1),\dots,(A_n,\alpha_n)\}$ un système tel que $S=\displaystyle\sum_{i=1}^{n}\alpha_i\neq0$. Le \textbf{barycentre} $G$ est l'unique point vérifiant
\[
\sum_{i=1}^{n}\alpha_i\,\overrightarrow{GA_i}=\vec 0
\qquad\Longleftrightarrow\qquad
\forall M,\ \sum_{i=1}^{n}\alpha_i\,\overrightarrow{MA_i}=S\,\overrightarrow{MG}.
\]
Lorsque tous les coefficients sont \textbf{égaux} (et non nuls), $G$ est l'\textbf{isobarycentre} des points $A_i$.
\end{definition}

\begin{remarque}
L'isobarycentre de deux points est leur \textbf{milieu} ; celui des trois sommets d'un triangle est son \textbf{centre de gravité} (intersection des médianes).
\end{remarque}

\begin{illus}
\begin{tikzpicture}[>=Stealth,scale=1.0]
\coordinate (A) at (0,0);
\coordinate (B) at (5,0.4);
\coordinate (C) at (1.6,3);
% G = barycentre de (A,1),(B,1),(C,2) : G=(A+B+2C)/4
\coordinate (G) at ($0.25*(A)+0.25*(B)+0.5*(C)$);
\draw[onbline] (A)--(B)--(C)--cycle;
\filldraw[onbink] (A) circle (1.6pt) node[below left]{\small $A\,(1)$};
\filldraw[onbink] (B) circle (1.6pt) node[below right]{\small $B\,(1)$};
\filldraw[onbink] (C) circle (1.6pt) node[above]{\small $C\,(2)$};
\filldraw[onbo] (G) circle (2pt) node[right]{\small $G$};
\draw[onbo,->,line width=0.9pt] (G)--(A);
\draw[onbo,->,line width=0.9pt] (G)--(B);
\draw[onbo,->,line width=0.9pt] (G)--(C);
\node[align=left,text width=4.5cm,anchor=west] at (6.1,1.5)
{\small $G=\mathrm{bar}\{(A,1),(B,1),(C,2)\}$ vérifie $\overrightarrow{GA}+\overrightarrow{GB}+2\,\overrightarrow{GC}=\vec0$. Le poids plus fort en $C$ « tire » $G$ vers $C$.};
\end{tikzpicture}
\fig{Figure — Barycentre de trois points pondérés.}
\end{illus}

\section{Propriétés fondamentales}

\begin{propriete}[Homogénéité]
Le barycentre est inchangé si l'on multiplie \textbf{tous} les coefficients par un même réel $k\neq0$ :
\[
\mathrm{bar}\{(A_i,\alpha_i)\}=\mathrm{bar}\{(A_i,k\alpha_i)\}.
\]
\end{propriete}

\begin{remarque}
Grâce à l'homogénéité, on peut toujours se ramener à des coefficients commodes (entiers, ou de somme $1$). Si $\sum\alpha_i=1$, alors $\overrightarrow{OG}=\sum\alpha_i\,\overrightarrow{OA_i}$ pour toute origine $O$.
\end{remarque}

\begin{theoreme}[Associativité — barycentres partiels]
On ne change pas le barycentre d'un système en remplaçant plusieurs points par leur \textbf{barycentre partiel}, affecté de la \textbf{somme} de leurs coefficients (à condition que cette somme soit non nulle). Par exemple, si $\alpha+\beta\neq0$ et $H=\mathrm{bar}\{(A,\alpha),(B,\beta)\}$, alors
\[
\mathrm{bar}\{(A,\alpha),(B,\beta),(C,\gamma)\}=\mathrm{bar}\{(H,\alpha+\beta),(C,\gamma)\}.
\]
\end{theoreme}

\begin{aretenir}
L'associativité est l'outil-roi pour \textbf{construire} un barycentre et pour prouver des \textbf{alignements} ou des \textbf{concours de droites} : regrouper deux points fait apparaître que $G$ est sur une droite remarquable.
\end{aretenir}

\begin{propriete}[Coordonnées du barycentre]
Dans un repère, si $A_i$ a pour coordonnées $(x_i,y_i)$ et $S=\sum\alpha_i\neq0$, alors
\[
x_G=\frac{1}{S}\sum_{i=1}^{n}\alpha_i x_i,\qquad
y_G=\frac{1}{S}\sum_{i=1}^{n}\alpha_i y_i.
\]
(De même pour l'affixe : $z_G=\frac1S\sum\alpha_i z_i$.)
\end{propriete}

\begin{exemple}
Barycentre de $(A,2),(B,1),(C,1)$ avec $A(1,0)$, $B(3,2)$, $C(0,-2)$ : $S=4$, donc
$x_G=\frac14(2\cdot1+3+0)=\frac54$, $y_G=\frac14(0+2-2)=0$. Ainsi $G\!\left(\frac54,0\right)$.
\end{exemple}

% ═══════════════════════════════════════════════════════════════
\section{Réduction de sommes vectorielles}

\begin{theoreme}[Réduction de $\sum\alpha_i\,\overrightarrow{MA_i}$]
Soit $\vec f(M)=\displaystyle\sum_{i=1}^{n}\alpha_i\,\overrightarrow{MA_i}$ et $S=\sum\alpha_i$.
\begin{itemize}
\item Si $S\neq0$ : soit $G$ le barycentre du système ; alors $\vec f(M)=S\,\overrightarrow{MG}$ (le vecteur dépend de $M$).
\item Si $S=0$ : $\vec f(M)$ est un vecteur \textbf{constant} (indépendant de $M$).
\end{itemize}
\end{theoreme}

\begin{exemple}
$\vec f(M)=2\,\overrightarrow{MA}-2\,\overrightarrow{MB}$ : ici $S=0$, donc $\vec f(M)=2(\overrightarrow{MA}-\overrightarrow{MB})=2\,\overrightarrow{BA}$, vecteur constant.
\end{exemple}

\section{Réduction de $\sum\alpha_i\,MA_i^{2}$ — formule de Leibniz}

\begin{theoreme}[Formule de Leibniz]
Posons $f(M)=\displaystyle\sum_{i=1}^{n}\alpha_i\,MA_i^{2}$ et $S=\sum\alpha_i$.
\begin{itemize}
\item \textbf{Si $S\neq0$}, soit $G$ le barycentre du système. Pour tout point $M$ :
\[
f(M)=S\cdot MG^{2}+f(G),\qquad\text{où } f(G)=\sum_{i=1}^{n}\alpha_i\,GA_i^{2}.
\]
\item \textbf{Si $S=0$}, soit $\vec u=\sum\alpha_i\,\overrightarrow{OA_i}$ (constant) ; alors $f(M)=2\,\overrightarrow{MO}\cdot\vec u+f(O)$ est une fonction \textbf{affine} de $M$.
\end{itemize}
\end{theoreme}

\begin{remarque}
Démonstration (cas $S\neq0$) : pour chaque $i$, $MA_i^{2}=\overrightarrow{MA_i}\cdot\overrightarrow{MA_i}=\big(\overrightarrow{MG}+\overrightarrow{GA_i}\big)^{2}=MG^{2}+2\,\overrightarrow{MG}\cdot\overrightarrow{GA_i}+GA_i^{2}$. En sommant pondéré par $\alpha_i$ : le terme du milieu vaut $2\,\overrightarrow{MG}\cdot\big(\sum\alpha_i\overrightarrow{GA_i}\big)=\vec0$ par définition de $G$. D'où $f(M)=S\,MG^{2}+f(G)$.
\end{remarque}

\begin{aretenir}
\textbf{Leibniz transforme} $\sum\alpha_i MA_i^{2}$ en $S\,MG^{2}+\text{cte}$. Une équation $\sum\alpha_i MA_i^{2}=k$ devient $MG^{2}=\dfrac{k-f(G)}{S}$ : c'est un \textbf{cercle} de centre $G$ (si le second membre $>0$), un point (si $=0$), ou l'ensemble vide (si $<0$).
\end{aretenir}

% ═══════════════════════════════════════════════════════════════
\section{Lignes de niveau}

\begin{definition}[Ligne de niveau]
Étant donnée une fonction $f$ qui à tout point $M$ associe un réel $f(M)$, la \textbf{ligne de niveau} $k$ est l'ensemble $\mathcal L_k=\{\,M\ :\ f(M)=k\,\}$.
\end{definition}

\begin{propriete}[Trois lignes de niveau classiques]
\begin{enumerate}[label=\arabic*)]
\item $f(M)=\sum\alpha_i MA_i^{2}$ : par Leibniz, $\mathcal L_k$ est un \textbf{cercle} de centre le barycentre $G$ (ou $\varnothing$, ou un point) si $S\neq0$, et une \textbf{droite} perpendiculaire à $\vec u$ si $S=0$.
\item $f(M)=\dfrac{MA}{MB}=k>0$ (avec $A\neq B$) : si $k=1$, c'est la \textbf{médiatrice} de $[AB]$ ; si $k\neq1$, c'est un \textbf{cercle} (cercle d'Apollonius) de diamètre $[IJ]$, où $I=\mathrm{bar}\{(A,1),(B,-k)\}$ et $J=\mathrm{bar}\{(A,1),(B,k)\}$.
\item $f(M)=\big(\overrightarrow{MA},\overrightarrow{MB}\big)=\theta$ (angle orienté, $M\neq A,B$) : c'est un \textbf{arc de cercle} passant par $A$ et $B$ (« arc capable » d'angle $\theta$). Le cas $\theta=\pm\frac\pi2$ donne le \textbf{cercle de diamètre} $[AB]$.
\end{enumerate}
\end{propriete}

\begin{illus}
\begin{tikzpicture}[>=Stealth,scale=1.0]
% Ligne de niveau MA^2+MB^2 = cte : cercle centré au milieu I de [AB]
\coordinate (A) at (-2,0);
\coordinate (B) at (2,0);
\coordinate (I) at (0,0);
\draw[onbblue,line width=1pt] (I) circle (1.7);
\filldraw[onbink] (A) circle (1.6pt) node[below left]{\small $A$};
\filldraw[onbink] (B) circle (1.6pt) node[below right]{\small $B$};
\filldraw[onbo] (I) circle (1.8pt) node[below]{\small $I$};
% un point M sur le cercle
\coordinate (M) at (60:1.7);
\filldraw[onbgreen] (M) circle (1.6pt) node[above right]{\small $M$};
\draw[onbgreen,->] (M)--(A);
\draw[onbgreen,->] (M)--(B);
\node[align=left,text width=4.6cm,anchor=west] at (2.6,0.4)
{\small $MA^{2}+MB^{2}=k$ : ici $S=2$, le barycentre est le milieu $I$ de $[AB]$. La ligne de niveau est le \textbf{cercle} de centre $I$ et de rayon $\sqrt{\tfrac{k}{2}-\tfrac{AB^{2}}{4}}$.};
\end{tikzpicture}
\fig{Figure — Ligne de niveau $M\mapsto MA^{2}+MB^{2}$ : un cercle centré au milieu de $[AB]$.}
\end{illus}

% ═══════════════════════════════════════════════════════════════
\section{Méthodes \& savoir-faire}

\methode{Réduire une somme vectorielle pondérée}
Calculer $S=\sum\alpha_i$. Si $S\neq0$, introduire le barycentre $G$ : $\sum\alpha_i\overrightarrow{MA_i}=S\,\overrightarrow{MG}$. Si $S=0$, regrouper pour faire apparaître un vecteur constant.
\begin{exemple}
Réduire $\vec f(M)=\overrightarrow{MA}+2\,\overrightarrow{MB}+\overrightarrow{MC}$. Ici $S=4\neq0$ ; soit $G=\mathrm{bar}\{(A,1),(B,2),(C,1)\}$. Alors $\vec f(M)=4\,\overrightarrow{MG}$.
\end{exemple}

\methode{Réduire $\sum\alpha_i\,MA_i^{2}$ par Leibniz}
Calculer $S$, le barycentre $G$, puis la constante $f(G)=\sum\alpha_i\,GA_i^{2}$. On obtient $f(M)=S\,MG^{2}+f(G)$.
\begin{exemple}
$A,B$ avec $AB=4$, $f(M)=MA^{2}+MB^{2}$. $S=2$, $G=I$ milieu de $[AB]$, $GA=GB=2$, donc $f(G)=2^{2}+2^{2}=8$ et $f(M)=2\,MI^{2}+8$.
\end{exemple}

\methode{Déterminer et construire une ligne de niveau}
Réduire $f(M)$ (vectoriel ou Leibniz), isoler $MG^{2}=R^{2}$ : reconnaître cercle, point ou $\varnothing$ ; si $S=0$, on obtient une droite.
\begin{exemple}
$AB=4$ : déterminer $\{M:MA^{2}+MB^{2}=24\}$. D'après ci-dessus $2\,MI^{2}+8=24$, soit $MI^{2}=8$, $MI=2\sqrt2$ : \textbf{cercle} de centre $I$ et de rayon $2\sqrt2$.
\end{exemple}

\methode{Construire un barycentre par associativité}
Regrouper deux points en un barycentre partiel $H$, puis placer $G$ sur $(HC)$ via $\overrightarrow{HG}=\frac{\gamma}{S}\overrightarrow{HC}$.
\begin{exemple}
$G=\mathrm{bar}\{(A,1),(B,1),(C,2)\}$ : poser $H=$ milieu de $[AB]$ ($\mathrm{bar}\{(A,1),(B,1)\}$), alors $G=\mathrm{bar}\{(H,2),(C,2)\}=$ milieu de $[HC]$.
\end{exemple}

\methode{Prouver un alignement ou un concours de droites}
Montrer qu'un point est barycentre de deux autres : il est alors \textbf{sur la droite} qui les joint. Pour un concours, exhiber un même barycentre $G$ appartenant à chaque droite.
\begin{exemple}
Centre de gravité d'un triangle : $G=\mathrm{bar}\{(A,1),(B,1),(C,1)\}$. En regroupant $B,C$ en leur milieu $A'$ : $G=\mathrm{bar}\{(A,1),(A',2)\}\in(AA')$. De même $G\in(BB')$ et $(CC')$ : les trois médianes \textbf{concourent} en $G$, avec $\overrightarrow{AG}=\frac23\overrightarrow{AA'}$.
\end{exemple}

\methode{Reconnaître une ligne de niveau $MA/MB$ ou un angle}
Pour $MA/MB=k$ : si $k=1$ médiatrice de $[AB]$, sinon cercle d'Apollonius. Pour un angle $(\overrightarrow{MA},\overrightarrow{MB})$ constant : arc capable ; angle droit $\Rightarrow$ cercle de diamètre $[AB]$.
\begin{exemple}
$\{M:MA=MB\}$ est la médiatrice de $[AB]$ ; $\{M:(\overrightarrow{MA},\overrightarrow{MB})=\frac\pi2\}$ est le cercle de diamètre $[AB]$ (privé de $A$ et $B$).
\end{exemple}

% ═══════════════════════════════════════════════════════════════
\section{Exercices}

\rubrique{Barycentres et coordonnées}
\exo{}\diff{1} On donne $A(2,1)$, $B(-1,3)$, $C(4,-2)$. Déterminer les coordonnées du barycentre $G$ de $\{(A,1),(B,2),(C,1)\}$.

\exo{}\diff{1} Soient $A$ et $B$ deux points. Construire le barycentre $G$ de $\{(A,3),(B,1)\}$ en exprimant $\overrightarrow{AG}$ en fonction de $\overrightarrow{AB}$.

\exo{}\diff{2} À l'aide de l'homogénéité et de l'associativité, construire $G=\mathrm{bar}\{(A,2),(B,1),(C,1)\}$ en regroupant $B$ et $C$.

\rubrique{Réductions vectorielles}
\exo{}\diff{1} Réduire $\vec f(M)=3\,\overrightarrow{MA}-\overrightarrow{MB}$ et $\vec g(M)=\overrightarrow{MA}+\overrightarrow{MB}-2\,\overrightarrow{MC}$.

\exo{}\diff{2} Soit $ABC$ un triangle. Déterminer l'ensemble des points $M$ tels que $\big\|\,2\,\overrightarrow{MA}+\overrightarrow{MB}+\overrightarrow{MC}\,\big\|=8$.

\rubrique{Formule de Leibniz et lignes de niveau}
\exo{}\diff{2} $A$ et $B$ tels que $AB=6$. Déterminer et construire l'ensemble $\{M:MA^{2}+MB^{2}=26\}$.

\exo{}\diff{2} Mêmes $A,B$ ($AB=6$). Déterminer l'ensemble $\{M:MA^{2}-MB^{2}=12\}$ (on remarquera que $S=0$).

\exo{}\diff{3} Soit $ABC$ un triangle équilatéral de côté $a$. Déterminer le lieu des points $M$ tels que $MA^{2}+MB^{2}+MC^{2}=2a^{2}$.

\rubrique{Lignes de niveau $MA/MB$ et angle}
\exo{}\diff{2} $A$ et $B$ distincts. Décrire l'ensemble $\{M:MA=2\,MB\}$ (cercle d'Apollonius).

\exo{}\diff{2} Décrire $\{M:(\overrightarrow{MA},\overrightarrow{MB})=\frac\pi2\}$ puis $\{M:MA=MB\}$.

\rubrique{Problème type Baccalauréat}
\exo{}\diff{3} Le plan est rapporté à un repère orthonormé. On donne $A(0,0)$, $B(4,0)$ et $C(0,3)$.
\begin{enumerate}[label=\arabic*)]
\item Déterminer les coordonnées du barycentre $G$ de $\{(A,2),(B,1),(C,1)\}$.
\item Réduire $\vec u(M)=2\,\overrightarrow{MA}+\overrightarrow{MB}+\overrightarrow{MC}$, puis déterminer l'ensemble $(\Gamma_1)$ des points $M$ tels que $\big\|\vec u(M)\big\|=8$.
\item On pose $f(M)=2\,MA^{2}+MB^{2}+MC^{2}$. Calculer $f(G)$, puis montrer que $f(M)=4\,MG^{2}+f(G)$.
\item Déterminer et tracer l'ensemble $(\Gamma_2)$ des points $M$ tels que $f(M)=49$.
\end{enumerate}

% ═══════════════════════════════════════════════════════════════
\section{Corrigés}

\corrige{de l'exercice 1}
$S=1+2+1=4$. $x_G=\frac14(2+2(-1)+4)=\frac14\cdot4=1$ ; $y_G=\frac14(1+2\cdot3+(-2))=\frac14\cdot5=\frac54$. Donc $G\!\left(1,\frac54\right)$.

\corrige{de l'exercice 2}
$S=3+1=4$, et $\overrightarrow{AG}=\frac{\beta}{S}\overrightarrow{AB}=\frac14\overrightarrow{AB}$ : $G$ est le point de $[AB]$ situé au quart à partir de $A$.

\corrige{de l'exercice 3}
Par homogénéité, $S=4$. Soit $H=\mathrm{bar}\{(B,1),(C,1)\}=$ milieu de $[BC]$. Par associativité, $G=\mathrm{bar}\{(A,2),(H,2)\}=$ milieu de $[AH]$. On a donc $\overrightarrow{AG}=\frac12\overrightarrow{AH}$.

\corrige{de l'exercice 4}
Pour $\vec f$ : $S=3-1=2\neq0$, soit $G=\mathrm{bar}\{(A,3),(B,-1)\}$ ; alors $\vec f(M)=2\,\overrightarrow{MG}$.
Pour $\vec g$ : $S=1+1-2=0$, donc $\vec g$ est constant : $\vec g(M)=\overrightarrow{MA}+\overrightarrow{MB}-2\,\overrightarrow{MC}=(\overrightarrow{MA}-\overrightarrow{MC})+(\overrightarrow{MB}-\overrightarrow{MC})=\overrightarrow{CA}+\overrightarrow{CB}$.

\corrige{de l'exercice 5}
$S=2+1+1=4\neq0$. Soit $G=\mathrm{bar}\{(A,2),(B,1),(C,1)\}$ ; alors $2\overrightarrow{MA}+\overrightarrow{MB}+\overrightarrow{MC}=4\,\overrightarrow{MG}$, donc la condition s'écrit $4\,MG=8$, soit $MG=2$. L'ensemble est le \textbf{cercle} de centre $G$ et de rayon $2$.

\corrige{de l'exercice 6}
$S=2$, $G=I$ milieu de $[AB]$, $IA=IB=3$, donc $f(I)=IA^{2}+IB^{2}=9+9=18$. Par Leibniz $f(M)=2\,MI^{2}+18=26$, d'où $MI^{2}=4$, $MI=2$ : \textbf{cercle} de centre $I$ (milieu de $[AB]$) et de rayon $2$.

\corrige{de l'exercice 7}
$S=1-1=0$. Posons $f(M)=MA^{2}-MB^{2}=\overrightarrow{MA}^{2}-\overrightarrow{MB}^{2}=(\overrightarrow{MA}-\overrightarrow{MB})\cdot(\overrightarrow{MA}+\overrightarrow{MB})=\overrightarrow{BA}\cdot(\overrightarrow{MA}+\overrightarrow{MB})$.
Avec $I$ milieu de $[AB]$ : $\overrightarrow{MA}+\overrightarrow{MB}=2\,\overrightarrow{MI}$, donc $f(M)=2\,\overrightarrow{BA}\cdot\overrightarrow{MI}=12$, soit $\overrightarrow{BA}\cdot\overrightarrow{MI}=6$. C'est une \textbf{droite perpendiculaire} à $(AB)$. (On a $\overrightarrow{BA}\cdot\overrightarrow{IM}=-6$ ; le pied $H$ vérifie $\overrightarrow{IH}=\frac{-6}{AB^{2}}\overrightarrow{BA}=\frac{-6}{36}\overrightarrow{BA}=-\frac16\overrightarrow{BA}=\frac16\overrightarrow{AB}$.) La droite coupe $(AB)$ au point $H$ tel que $\overrightarrow{IH}=\frac16\overrightarrow{AB}$, du côté de $A$.

\corrige{de l'exercice 8}
$S=3\neq0$, le barycentre est le centre de gravité $G$ du triangle. Calculons $f(G)=GA^{2}+GB^{2}+GC^{2}$. Pour un triangle équilatéral de côté $a$, $GA=GB=GC=\frac{a}{\sqrt3}$ (rayon du cercle circonscrit), donc $f(G)=3\cdot\frac{a^{2}}{3}=a^{2}$. Par Leibniz $f(M)=3\,MG^{2}+a^{2}=2a^{2}$, d'où $MG^{2}=\frac{a^{2}}{3}$, $MG=\frac{a}{\sqrt3}$ : \textbf{cercle} de centre $G$ et de rayon $\frac{a}{\sqrt3}$ — c'est exactement le \textbf{cercle circonscrit} au triangle.

\corrige{de l'exercice 9}
$MA=2MB\Leftrightarrow MA^{2}=4MB^{2}\Leftrightarrow MA^{2}-4MB^{2}=0$. Ici $S=1-4=-3\neq0$ : soit $G=\mathrm{bar}\{(A,1),(B,-4)\}$. Par Leibniz, $MA^{2}-4MB^{2}=-3\,MG^{2}+\big(GA^{2}-4GB^{2}\big)$. La condition équivaut à $MG^{2}=\frac{GA^{2}-4GB^{2}}{3}=$ constante $>0$ : c'est un \textbf{cercle} de centre $G$ (cercle d'Apollonius). Ses points sur $(AB)$ sont $I=\mathrm{bar}\{(A,1),(B,-2)\}$ et $J=\mathrm{bar}\{(A,1),(B,2)\}$, qui en forment un diamètre.

\corrige{de l'exercice 10}
$\{M:(\overrightarrow{MA},\overrightarrow{MB})=\frac\pi2\}$ : l'angle $\widehat{AMB}$ est droit, donc $M$ voit $[AB]$ sous un angle droit : c'est le \textbf{cercle de diamètre} $[AB]$ (privé de $A$ et $B$).
$\{M:MA=MB\}$ : ensemble des points équidistants de $A$ et $B$, c'est la \textbf{médiatrice} de $[AB]$.

\corrige{du problème}
\textbf{1)} $S=2+1+1=4$. $x_G=\frac14(2\cdot0+4+0)=1$ ; $y_G=\frac14(0+0+3)=\frac34$. Donc $G\!\left(1,\frac34\right)$.

\textbf{2)} $S=4\neq0$ donc $\vec u(M)=2\overrightarrow{MA}+\overrightarrow{MB}+\overrightarrow{MC}=4\,\overrightarrow{MG}$. La condition $\|\vec u(M)\|=8$ s'écrit $4\,MG=8$, soit $MG=2$ : $(\Gamma_1)$ est le \textbf{cercle} de centre $G\!\left(1,\frac34\right)$ et de rayon $2$.

\textbf{3)} Calculons $f(G)=2\,GA^{2}+GB^{2}+GC^{2}$.
$GA^{2}=(1-0)^{2}+\left(\frac34\right)^{2}=1+\frac{9}{16}=\frac{25}{16}$.
$GB^{2}=(1-4)^{2}+\left(\frac34\right)^{2}=9+\frac{9}{16}=\frac{153}{16}$.
$GC^{2}=(1-0)^{2}+\left(\frac34-3\right)^{2}=1+\frac{81}{16}=\frac{97}{16}$.
Donc $f(G)=2\cdot\frac{25}{16}+\frac{153}{16}+\frac{97}{16}=\frac{50+153+97}{16}=\frac{300}{16}=\frac{75}{4}$.
La formule de Leibniz (avec $S=4$) donne $f(M)=4\,MG^{2}+f(G)=4\,MG^{2}+\frac{75}{4}$.

\textbf{4)} $f(M)=49\Leftrightarrow 4\,MG^{2}+\frac{75}{4}=49\Leftrightarrow 4\,MG^{2}=49-\frac{75}{4}=\frac{196-75}{4}=\frac{121}{4}\Leftrightarrow MG^{2}=\frac{121}{16}$, soit $MG=\frac{11}{4}=2{,}75$. $(\Gamma_2)$ est le \textbf{cercle} de centre $G\!\left(1,\frac34\right)$ et de rayon $\frac{11}{4}$ (concentrique à $(\Gamma_1)$).
